\subsubsection{Backend}
El backend está implementado en el lenguaje de programación Rust, haciendo uso
del \textit{framework} Rocket \cite{rocket}. Consta de un servidor web con una 
API tipo REST, la cual se basa en el protocolo HTTP. Para la integración con
la base de datos de MongoDB, se usó la librería de Rust que proporciona el
propio Mongo.

El servicio cuenta con dos rutas para la interacción con los dos tipos de datos
implementados: \texttt{/sightings} y \texttt{/events}. Cada una de estas rutas
permite tres tipos de acciones:
\begin{itemize}
  \item Petición \texttt{GET} a \texttt{/<tipo>}: Retorna una lista de todos los
  elementos almacenados en base de datos correspondientes a ese tipo, realizando
  la correspondiente consulta.
  \item Petición \texttt{GET} a \texttt{/<tipo>/<id>}: Retorna los detalles del
  elemento requerido.
  \item Petición \texttt{POST} a \texttt{/<tipo>}: Almacena en base de datos el
  elemento especificado, siguiendo el esquema definido en la
  \subsecref{data-warehouse}.
\end{itemize}

En caso de que al devolver los elementos, éstos no cuenten con información
sobre la localización geográfica precisa del elemento, realiza una llamada a la
API de Mapbox \cite{mapbox-api} para obtenerla.

